\emph{Pour chacune des affirmations suivantes, indiquer si elle est vraie ou fausse, en justifiant la réponse.}

\begin{enumerate}
	\item On considère la suite arithmétique $\left(u_n\right)$ de raison $r=5$ et telle que $u_5=10$.

	\textbf{Affirmation :} $u_7 = 250$.

	\item Soit $\left(u_n\right)$ la suite définie pour tout entier naturel $n$ par $u_n = 2n^2+1$.

	\textbf{Affirmation :} la suite $\left(u_n\right)$ est une suite géométrique.

	\item Soit $f$ la fonction définie pour tout nombre réel $x$ par $f(x) = 5x^2-2x+3$.

	\textbf{Affirmation :} une équation de la tangente à la courbe représentative de $f$ au point d'abscisse 1 est donnée par :
	\[y = 8x-2.\]

	\item Avec la chaleur et la surpêche, une association de pêcheurs constate chaque année que le nombre de truites présentes dans l'étang de leur commune diminue de 15\,\%. À partir de 2024, pour enrayer ce phénomène, ils décident de réintroduire chaque année 100 truites dans l'étang. Pour tout entier naturel $n$, $u_n$ désigne le nombre de truites en $2024+n$.

	\textbf{Affirmation :} pour tout entier naturel $n$, $u_{n+1} = 0,15\,u_n + 100$.
\end{enumerate}
